%%%%%%%%%%%%%%%%%%%%%%%%%%%%%%%%%%%%%%%%%%%%%%%%%%%%%%%%%%%%%%%%%%%%%%%%%%%%%%
% engrain -- MLSys 2027 submission
%
% Every number in this paper is produced by a harness in this repository and
% names its artifact in the caption. Numbers that are stale with respect to the
% current implementation are not quoted at all; see paper/PROVENANCE.md.
%%%%%%%%%%%%%%%%%%%%%%%%%%%%%%%%%%%%%%%%%%%%%%%%%%%%%%%%%%%%%%%%%%%%%%%%%%%%%%

\documentclass{article}

\usepackage{microtype}
\usepackage{graphicx}
\usepackage{subfigure}
\usepackage{booktabs}
\usepackage{amsmath}
\usepackage{amssymb}
\usepackage{multirow}
\usepackage{hyperref}

\newcommand{\theHalgorithm}{\arabic{algorithm}}

\usepackage{mlsys2025}

\mlsystitlerunning{Engrain: Device-Resident LR(1) Constrained Decoding}

\begin{document}

\twocolumn[
\mlsystitle{Engrain: Constrained Decoding with the Parser Resident on the GPU}

\begin{mlsysauthorlist}
\mlsysauthor{Anonymous Authors}{anon}
\end{mlsysauthorlist}

\mlsysaffiliation{anon}{Anonymous Institution}

\mlsyskeywords{constrained decoding, structured generation, LR parsing, CUDA
graphs, LLM serving}

\vskip 0.3in

\begin{abstract}
Constrained decoding restricts a language model to a formal language by masking
the logits at every step, and every deployed implementation computes that mask
on the host. This is not an efficiency detail: a CUDA graph is a recorded
sequence of \emph{device} work, so a component that requires the host mid-step
cannot be inside the graph a serving engine replays for its decode loop, cannot
hand an allowed set to a fused sampler, and cannot follow a speculative draft
without a synchronization per position. We present \textbf{engrain}, a
constrained-decoding engine whose LALR(1) parser state --- lexer state, parser
stack, and conflict forks --- lives in GPU memory and advances in device
kernels, so that the grammar is part of the decode step rather than an input to
it. The design contributes an exact device encoding of the (lexer
state $\times$ parser stack) configuration space, a mask fill that is a single
launch for an arbitrarily heterogeneous batch, and a cross-step memo that
answers a third of all steps without touching the tables. On an A100 with a
151{,}669-token vocabulary, a full constrained step (mask \emph{and} parser
advance) costs 103--327\,\textmu{}s at batch 512 against XGrammar's
2{,}041--6{,}811\,\textmu{}s, adds 90\,\textmu{}s to an overlapped decode step
where a host matcher adds 398, holds its p99 flat when the host loses cores to
other work while the host matcher's p99 rises 7.8$\times$, and is
1.35--4.22$\times$ cheaper under speculative decoding. Masks are bit-exact
against a reference matcher over 7{,}198 verified rows. We also report where
residency is paid for: at batch 1 we are 0.24--0.90$\times$, our compiler is
7$\times$ slower, our resident tables are 27$\times$ larger at the median, and
our lowering still accepts documents its schema rejects 19.6\% of the time
against XGrammar's 5.4\% --- an audit that located and fixed a silent
\texttt{required}-dropping defect along the way.
\end{abstract}
]

\printAffiliationsAndNotice{}

\section{Introduction}
\label{sec:intro}

Structured generation --- forcing a language model's output to be valid JSON,
to match a schema, or to parse under a grammar --- is now a standard serving
feature. It is implemented by \emph{constrained decoding}: before each token is
sampled, the engine computes a bitmask over the vocabulary marking which tokens
can extend the output without leaving the language, and sets the rest to
$-\infty$.

Every widely deployed implementation --- XGrammar~\cite{dong2024xgrammar},
llguidance, Outlines~\cite{willard2023efficient}, SynCode~\cite{ugare2025syncode}
--- computes that mask with a CPU automaton and copies it to the device. The
choice looks reasonable: the mask is small, the automaton step is a few
microseconds, and the copy overlaps the forward pass. Measured as a fraction of
a decode step, it is 4.6\% at batch 512 on a 0.6B model, and an engine that
made it twice as fast would be 2.3\% faster end to end. Treated as a latency
problem, constrained decoding is nearly finished.

This paper argues that latency is the wrong frame, and that the interesting
property is \emph{where the state lives}.

Modern serving engines record their decode step as a CUDA graph and replay it,
because at batch 512 the per-kernel launch cost otherwise dominates. A CUDA
graph is a fixed, recorded sequence of device work. Anything the host has to
\emph{produce} in the middle of a step cannot be recorded into it --- attempting
to capture a host-side mask fill yields a graph that contains none of the work,
and replaying it reproduces whatever the host buffer happened to hold last.
Consequently, every deployed constrained decoder sits \emph{outside} the graph
and hands a buffer in.

Three things follow, and none of them is a speed argument:

\begin{itemize}
\item \textbf{The decode loop cannot be captured whole.} A host-dependent
component cannot join a fully graph-resident or megakernel decode loop at all.
\item \textbf{The sampler cannot be fused with the constraint.} If the allowed
set is a device object, a sampler can read the few hundred candidates the
grammar permits instead of all 151{,}669 logits. If it is produced on the host,
the only interface is a vocabulary-wide mask.
\item \textbf{Speculative decoding needs a host round trip per draft position},
or a specialized batch API, because the parser cannot be advanced and rolled
back inside the same recorded step.
\end{itemize}

On each of these the comparison is not ``we are faster'' but ``the other design
cannot participate''. That is the claim this paper makes, and the engineering it
requires is the paper's substance: moving an LR parser onto a GPU is not a port.

\paragraph{Why an LR parser, and why that is hard.} A grammar-constrained
decoder must answer, at every step, ``which of the model's tokens may come
next''. Model tokens are byte strings from a BPE vocabulary; they straddle
lexeme boundaries, so a token may end in the middle of a number or a string. The
parser configuration must therefore carry a \emph{lexer} state as well as a
parser stack. Worse, an LR parser's next state after a reduction depends on the
stack below the popped states, so \texttt{next[state, token]} is not a function
of the state alone and no flat per-state table is exact. Prior GPU work avoids
this by restricting the grammar class: an acyclic schema DFA, a bounded stack
depth, or a conservative stack abstraction. Restricting the class is precisely
what makes those systems unable to serve the grammars people write.

\paragraph{Contributions.}
\begin{enumerate}
\item An exact device representation of the (lexer state $\times$ parser stack)
configuration space for LALR(1) grammars, including runtime forking on
reduce/reduce conflicts, so that no bound truncates the accepted language
(\S\ref{sec:automaton}).
\item A mask fill that is a single kernel launch for a batch of arbitrarily
many distinct grammars at arbitrarily different document positions, built as a
\emph{probe} that discovers per-sequence work and a \emph{sweep} that
redistributes it --- and evidence that the obvious single-kernel fusion is
16.4$\times$ slower on the batch shape serving actually produces
(\S\ref{sec:fill}).
\item A cross-step memo, keyed on the configuration set, that answers 35.6\% of
real decoding steps without reading the tables (\S\ref{sec:memo}).
\item Fused constrained sampling: the allowed set is compacted on device and
consumed directly by the sampler (\S\ref{sec:sampling}), together with an
honest account of why the win does not survive a ragged batch.
\item A differential correctness methodology, and 7{,}198 verified mask rows
across four independent verifications (\S\ref{sec:correctness}).
\item An evaluation that reports the costs of residency --- compile time,
resident memory, small-batch latency --- as first-class results rather than
limitations (\S\ref{sec:costs}).
\end{enumerate}

\section{Background and Motivation}
\label{sec:background}

\subsection{Constrained decoding}

Given a context-free grammar $G$ over a terminal alphabet and a model vocabulary
$V$ of byte strings, a constrained decoder maintains, per sequence, an automaton
state and computes at each step the set
$A \subseteq V$ of tokens $t$ such that the output so far concatenated with $t$
is a prefix of some string in $L(G)$. The mask is the indicator of $A$.

XGrammar~\cite{dong2024xgrammar} --- the most widely deployed implementation,
and our baseline throughout --- compiles the grammar to a pushdown automaton,
precomputes for each automaton state an ``adaptive token mask'' partitioned into
accepted and rejected sets, and at each step walks the per-sequence stack on the
CPU to fill a bitmask, which is copied to the device and applied by a small
kernel. Its compiler cache is kilobyte-scale and its compile time is
milliseconds, both of which are properties we do not match and report honestly
in \S\ref{sec:costs}.

\subsection{The graph boundary}

At batch 512, a decode step of a small model issues hundreds of kernels whose
individual launch overheads exceed their execution times. Serving engines
therefore record the step once as a CUDA graph and replay it. Capture is
\emph{stream capture}: the operations enqueued on a stream between
\texttt{cudaStreamBeginCapture} and \texttt{cudaStreamEndCapture} become graph
nodes. Host code executes normally during capture and contributes nothing.

This makes the boundary sharp rather than gradual. A constrained decoder that
computes its mask on the host is not ``slower inside a graph''; it is absent
from the graph, and the replayed step silently uses a stale buffer. The only
sound arrangement is the one deployed today: keep the fill outside, synchronize,
hand in a buffer. Every property in \S\ref{sec:intro} follows from that
arrangement, not from the fill's cost.

\subsection{What the workload actually looks like}

Two properties of real constrained decoding shape the design, and both are
measured over model-generated documents rather than assumed. We generate 533
JSON values with Llama-3-8B-Instruct under an XGrammar constraint (50 schemas
from each of the 11 JSONSchemaBench~\cite{geng2025jsonschemabench} configs,
temperature 0.8, top-$p$ 0.95, mean 125 tokens) and replay the text through
each tokenizer's grammar, which is faithful because the state a matcher reaches
is determined by the bytes consumed.

\begin{table}[t]
\caption{Allowed-set width over real decoding steps. The median step allows a
few hundred tokens regardless of vocabulary size, but a third of steps allow
more than 8{,}192 --- the distribution is bimodal, not concentrated.
Source: \texttt{results/rigorous-summary.json}.}
\label{tab:width}
\vskip 0.1in
\begin{center}
\begin{small}
\begin{tabular}{lrrrr}
\toprule
tokenizer & $|V|$ & steps & median $|A|$ & $|A|>8$k \\
\midrule
Llama 3      & 128{,}256 & 55{,}406 & 396 & 32.9\% \\
Qwen 3.6     & 248{,}077 & 66{,}557 & 378 & 32.0\% \\
Gemma 4      & 262{,}144 & 69{,}313 & 107 & 32.3\% \\
\bottomrule
\end{tabular}
\end{small}
\end{center}
\vskip -0.1in
\end{table}

Table~\ref{tab:width} shows the first property: \textbf{width is set by the
schema, not the vocabulary}. A structural position (after a comma, inside an
object) allows a few hundred tokens whatever the vocabulary size, so the ratio
between an $O(|V|)$ mask and an $O(|A|)$ allowed set \emph{grows} as vocabularies
grow. The second property is that the distribution is bimodal: string bodies
allow nearly everything, and a third of all steps are in one. Any design that
assumes narrow rows --- including our own fused sampler --- meets that third of
steps, and \S\ref{sec:sampling} reports what happens when it does.

The third property is about batch shape, and it caused the largest single
regression in this project. A serving batch is \emph{skewed by construction}:
continuous batching admits requests at different times, so every sequence sits
at its own position in its own document, and the amount of parser work each
needs differs by an order of magnitude. Benchmarks that place every sequence in
the same parse state measure the one arrangement that flatters a design which
deduplicates. All measurements in this paper seed each sequence at a random
position in its document.

\section{Design}
\label{sec:design}

\subsection{The automaton, and where it ends}
\label{sec:automaton}

A configuration is a pair (lexer DFA state, parser stack). The compiler produces,
from a JSON Schema, EBNF or regex front end: an LALR(1) ACTION/GOTO table over
grammar terminals; a byte-level lexer DFA whose accepting states name terminals;
and a partition of the model vocabulary into \emph{token groups}.

\paragraph{Token groups.} The bridge between BPE tokens and grammar terminals is
the expensive object. For each lexer state $\ell$ and each token $t$, feeding
$t$'s bytes from $\ell$ either fails, or ends in some lexer state $\ell'$ having
emitted a (possibly empty) sequence of terminals. Two tokens that do the same
thing from \emph{every} lexer state are interchangeable for masking purposes, so
the vocabulary is quotiented by that equivalence. Real vocabularies collapse
hard: 151{,}669 Qwen3 tokens become a few hundred groups per grammar. The mask
fill then works over groups, and expands to tokens only at the end.

\paragraph{Stacks, exactly.} The classical obstacle is that a reduction pops
$k$ states and the following GOTO depends on what is exposed underneath, so a
flat \texttt{next[state, token]} table is not exact. Bounding the stack depth
and enumerating reachable stacks --- the approach a prototype of this system
took --- both explodes (compile timed out beyond depth 6 at a 4{,}096-token
vocabulary) and silently makes the accepted language a strict subset of the
grammar when the bound is exceeded. That is a correctness cliff, not a
performance one.

Engrain keeps the stack. A sequence's device state is a bounded-capacity array
of configurations, each holding a lexer state and an explicit parser stack, and
the advance kernel performs real reductions with real GOTO lookups. What is
bounded is the number of \emph{simultaneous} configurations, not the depth of
any stack, and exceeding that bound is reported as a typed refusal at compile
time rather than silently narrowing the language.

\paragraph{Conflicts, forked at runtime.} LALR(1) construction on lowered JSON
Schema produces reduce/reduce conflicts --- \texttt{oneOf} in 86\% of the
conflicted grammars. Rather than refuse those grammars (which would cost a fifth
of the corpus) or promote the whole system to full GLR, engrain forks the
conflicted configuration at runtime and carries both, pruning any fork that dies
on the next terminal. This is GLR-lite: unbounded in principle, bounded by the
configuration capacity in practice, and the capacity is checked rather than
assumed.

\subsection{The fill: a probe and a sweep}
\label{sec:fill}

The mask fill is the step's device-side work. Its input is, per sequence, a set
of configurations; its output is a bitmask row over the vocabulary. Conceptually
it is: for each sequence, for each configuration, for each token group reachable
from that configuration's lexer state, if the group's terminal is accepted by
ACTION in the configuration's parser state, OR the group's token bits into the
row.

The iteration space is therefore ragged in three dimensions at once ---
sequences have different numbers of configurations, configurations have
different numbers of groups, and groups have different numbers of tokens. Our
first CUDA implementation fused the whole thing into one kernel with one block
per sequence. On a balanced batch it was 2.3--4.5$\times$ faster than the Triton
implementation it replaced. On a \emph{skewed} batch --- the shape serving
actually produces --- it was \textbf{16.4$\times$ slower} (9{,}883\,\textmu{}s
against 604\,\textmu{}s), because one block per sequence makes the fill cost the
\emph{widest} sequence rather than the total work. Measuring a fused kernel on a
balanced batch is not measuring it.

The fix is two kernels, and the reason it must be two is a hazard rather than a
preference:

\begin{itemize}
\item \textbf{Probe} (one block per sequence) clears the sequence's row and
records how many (configuration, group) items it will contribute.
\item \textbf{Sweep} (grid $(\text{batch}, \text{chunks})$) walks a
\emph{flattened} (configuration, group) iteration space, so a sequence with 64
configurations of one group each is spread over blocks rather than serialized
into one.
\end{itemize}

They cannot be merged: the probe clears rows that the sweep ORs into, and a
block clearing a row while another accumulates into it loses bits. The chunk
count is bounded by a scratch budget rather than by the batch --- a thread owns
two replay windows, so the buffer is $\text{batch} \times \text{chunks} \times
\text{threads}$, which is 268\,MB at batch 512 with 8 chunks. A 64\,MiB budget
gives 32 chunks at batch 32 and 2 at batch 512.

After the split the skewed step fell from 9{,}883\,\textmu{}s to
641\,\textmu{}s, a 15.4$\times$ improvement, and the remaining gap to our own
Triton reference narrowed from 1.16--3.29$\times$ to 1.16--1.31$\times$.

\subsection{The advance}

After a token is sampled, each sequence's configurations must be advanced by it.
This is the half that cannot be hidden behind the forward pass, because it
follows the sampled token. Engrain fuses the whole advance --- feeding the
token's bytes through the lexer, running reductions and GOTOs, forking on
conflicts, compacting dead configurations, and writing the rollback history
entry --- into a single kernel. The rollback path went from 16--20 graph nodes
to 7 by folding the history write inline, removing a counting kernel and a
$\text{grid}=(1,)$ prefix sum.

\subsection{Cross-step memoisation}
\label{sec:memo}

Sequences repeat states. Within one request, consecutive steps often reach a
configuration set that has been seen before; across a batch under one grammar,
many sequences share one. Engrain hashes the configuration set on device --- one
warp per configuration, folding stacks in shared memory --- and looks it up in a
device-resident memo of previously computed rows. Over a real walk the memo
answers \textbf{35.6\%} of steps without reading a table.

The hash itself was a bottleneck until it was parallelized: one warp performing
64 sequential stack reductions cost 65.8\,\textmu{}s; one warp per configuration
costs 16.3\,\textmu{}s, half of the Triton reference's 32.8.

\subsection{Fused constrained sampling}
\label{sec:sampling}

Because the allowed set is a device object, the sampler can read it directly.
A \texttt{compact} kernel emits, per row, a sorted list of allowed token ids
using a block-wide scan. The sampler then gathers a few hundred logits instead
of reading all 151{,}669.

The honest accounting of this idea is in \S\ref{sec:eval-sampling}, and it is
not the one this paper originally claimed. Compaction wins by 7.7$\times$ on a
sparse row and \emph{loses} by 0.95$\times$ on a dense one, so a fixed policy is
a net loss on the real bimodal distribution; a per-step choice made on device
(emit whichever of the allowed set and its complement is shorter) recovers
1.79$\times$; and on a genuinely ragged batch --- half the rows dense, which is
the measured steady state --- the advantage falls to 1.05$\times$. The reason is
not in the parser: it is that XGrammar's mask-apply takes its row indices as a
host \texttt{List[int]}, and a sampler's output shape is its row count, so any
subset scheme needs a host synchronization that the residency claim exists to
avoid.

\subsection{Arena and paging}

All grammars the engine has seen live in one device arena, and a sequence
carries the index of the grammar it is under, which is what makes a batch of
distinct grammars a single launch. Under continuous batching the arena must
reclaim memory, and the hard part is doing so \emph{without moving anything}:
compaction renumbers survivors, and every recorded CUDA graph would have to be
re-recorded, giving back exactly what residency bought. Engrain therefore gives
each array a coalescing free list and recycles identifiers rather than
renumbering them, pinning any grammar a live sequence is running under and
evicting the rest least-recently-used. Across 400 admissions of 40 schemas at a
32\,MB budget, 399 evictions cost \textbf{zero} graph re-records.

\section{Implementation}
\label{sec:impl}

The compiler is 7 Rust crates ($\approx$\,20k lines): front ends (JSON Schema,
EBNF, regex), lexer construction and vocabulary grouping, LALR(1) table
construction, artifact emission, and a reference matcher. The device runtime is
CUDA compiled to a fatbin for sm\_80--100 plus forward-compatible PTX, embedded
in the extension so that no driver JIT happens at load; the Python layer is
$\approx$\,5k lines over PyO3 bindings.

The system keeps two complete backends: the CUDA kernels, and a Triton
implementation of the same algorithms that is retained as an executable
reference. A third \texttt{differential} backend runs both on every call and
compares, which is how kernel-level divergence is caught. All results in this
paper are from the CUDA backend unless stated.

\section{Experimental Design}
\label{sec:design-exp}

A systems paper's claims are only as good as the angles from which someone
tried to break them. We state each claim with the measurement that would
falsify it, and report all of them --- including the two that came out against
us and the one that found a bug in our own compiler.

\begin{table}[t]
\caption{The claims, and the measurement designed to falsify each. Every row is
reported in \S\ref{sec:eval}, including those whose outcome is negative.}
\label{tab:design}
\vskip 0.1in
\begin{center}
\begin{small}
\begin{tabular}{p{0.29\columnwidth}p{0.60\columnwidth}}
\toprule
claim & what would falsify it \\
\midrule
the fill can live in the decode graph &
a graph recorded on one grammar assignment replaying wrongly on another \\
residency is worth having &
XGrammar's fill overlapping the forward pass so completely that its added cost
equals ours \\
the host is a liability &
the host matcher's tail staying flat as CPU cores are taken away \\
we are cheaper per step &
llguidance or a better-tuned XGrammar closing the gap \\
each mechanism earns its place &
removing the memo, or capture, changing nothing \\
heterogeneity is affordable &
cost growing with the number of distinct schemas in the batch \\
masks are exact &
a device row differing from the reference matcher \\
the grammar is the schema &
documents the grammar accepts that the schema rejects \\
\bottomrule
\end{tabular}
\end{small}
\end{center}
\vskip -0.1in
\end{table}

Two design rules apply throughout. \textbf{Both halves are charged}: a step is a
mask fill \emph{and} a parser advance, and reporting only the fill flatters
whichever engine has the cheaper one. \textbf{Batches are skewed}: every
sequence is seeded at a random position in its own document, because a balanced
batch is the single arrangement that flatters a design which deduplicates, and
measuring an ablation on one would understate every mechanism at once.

\section{Evaluation}
\label{sec:eval}

\subsection{Setup}

All measurements are on one NVIDIA A100 80GB PCIe with Qwen3-0.6B's
151{,}669-token vocabulary unless stated, and use three JSONSchemaBench schemas
of increasing difficulty (schema 0, 1, 2). Both engines are charged for
\emph{both} halves of a step --- the mask fill and the parser advance --- because
reporting only the fill flatters whichever engine has the cheaper one.
XGrammar's thread count is swept over $\{1,2,4,8,16,\text{auto}\}$ and its best
is reported; its device work is CUDA-graphed, as ours is. Every sequence is
seeded at a random position in its document. We report medians of 100--200
repeats after warm-up.

\subsection{Step latency}
\label{sec:eval-step}

Table~\ref{tab:step} is the paper's central performance result and also its
clearest limitation. At batch 1 we lose on all three schemas; the crossover is
between batch 8 and 32 depending on schema; above it the gap widens
monotonically to 12.8--20.8$\times$ at batch 512.

\begin{table}[t]
\caption{Full constrained step (fill + advance), median \textmu{}s, three
schemas. ``Triton'' is our own reference implementation of the same algorithms,
included so the CUDA work is measured against the engine it replaces and not
only against the baseline. Source: \texttt{engrain\_lab.rigor.latency}.}
\label{tab:step}
\vskip 0.1in
\begin{center}
\begin{small}
\begin{tabular}{rrrrr}
\toprule
batch & schema & engrain & Triton & XGrammar \\
\midrule
\multirow{3}{*}{1}
 & 0 & 49 & --- & 12 \\
 & 1 & 46 & --- & 14 \\
 & 2 & 41 & --- & 37 \\
\midrule
\multirow{3}{*}{16}
 & 0 & 57 & --- & 70 \\
 & 1 & 133 & --- & 82 \\
 & 2 & 162 & --- & 363 \\
\midrule
\multirow{3}{*}{32}
 & 0 & 59 & 83 & 130 \\
 & 1 & 136 & 116 & 150 \\
 & 2 & 194 & 147 & 659 \\
\midrule
\multirow{3}{*}{128}
 & 0 & 63 & 123 & 537 \\
 & 1 & 139 & 156 & 607 \\
 & 2 & 257 & 204 & 2{,}073 \\
\midrule
\multirow{3}{*}{512}
 & 0 & 103 & 156 & 2{,}041 \\
 & 1 & 185 & 203 & 2{,}375 \\
 & 2 & 327 & 247 & 6{,}811 \\
\bottomrule
\end{tabular}
\end{small}
\end{center}
\vskip -0.1in
\end{table}

Two readings matter. First, the shape is structural rather than incidental:
XGrammar's cost is per sequence and scales with the batch, ours is a fixed
number of launches whose work scales with \emph{distinct} parse states, so the
gap is a property of the two designs and not of tuning. Second, the small-batch
loss is also structural: a step sweeps every live configuration's groups
whatever the batch, so our floor is 41--49\,\textmu{}s at batch 1 while
XGrammar's per-sequence walk costs 12--37\,\textmu{}s. We do not claim a
small-batch win, and a serving deployment that runs at batch 1 should not use
this system.

Table~\ref{tab:step} also shows that our CUDA implementation still trails our
own Triton reference on the hardest schema (1.24--1.32$\times$). The cause is
per-item throughput in the sweep, not load imbalance, and we report it because
an implementation that loses to its own reference has not finished.

\subsection{Heterogeneous batches}
\label{sec:eval-mixed}

The case the design exists for is a batch of \emph{different} grammars, which is
what continuous batching produces and what should hurt us: our fill deduplicates
rows that share a grammar, a lexer state and a stack, and a mixture has fewer
such rows to share, while our buffers are sized by maxima over the whole pool.
XGrammar's per-sequence host call is schema-agnostic, so mixing costs it nothing
structurally.

\begin{table}[t]
\caption{Cost of a batch drawn from $n$ distinct schemas, median \textmu{}s, and
the cost of mixing itself: the ratio of the mixed batch to the mean of the same
schemas run alone at the same batch size.
Source: \texttt{probe\_mixed\_cost}, 2026-07-31.}
\label{tab:mixed}
\vskip 0.1in
\begin{center}
\begin{small}
\begin{tabular}{rrrrrr}
\toprule
batch & $n$ & engrain & XGrammar & ratio & mix costs us \\
\midrule
\multirow{5}{*}{32}
 & 1  & 82    & 166    & 2.01$\times$ & 0.99$\times$ \\
 & 2  & 388   & 181    & 0.47$\times$ & 5.00$\times$ \\
 & 4  & 401   & 398    & 0.99$\times$ & 5.38$\times$ \\
 & 8  & 442   & 1{,}991 & 4.51$\times$ & 5.50$\times$ \\
 & 16 & 394   & 1{,}086 & 2.76$\times$ & 5.19$\times$ \\
\midrule
\multirow{5}{*}{128}
 & 1  & 87    & 529    & 6.09$\times$ & 1.00$\times$ \\
 & 2  & 390   & 663    & 1.70$\times$ & 4.74$\times$ \\
 & 4  & 406   & 1{,}595 & 3.92$\times$ & 5.11$\times$ \\
 & 8  & 1{,}614 & 9{,}594 & 5.94$\times$ & 18.60$\times$ \\
 & 16 & 453   & 7{,}015 & 15.50$\times$ & 5.54$\times$ \\
\midrule
\multirow{5}{*}{512}
 & 1  & 114   & 1{,}930 & 16.89$\times$ & 0.99$\times$ \\
 & 2  & 419   & 2{,}456 & 5.86$\times$ & 3.79$\times$ \\
 & 4  & 430   & 6{,}604 & 15.37$\times$ & 4.02$\times$ \\
 & 8  & 4{,}025 & 35{,}781 & 8.89$\times$ & 33.17$\times$ \\
 & 16 & 4{,}052 & 44{,}963 & 11.10$\times$ & 35.98$\times$ \\
\bottomrule
\end{tabular}
\end{small}
\end{center}
\vskip -0.1in
\end{table}

Table~\ref{tab:mixed} reports both the comparison and the control. Three
findings, one of them against us:

\textbf{Heterogeneity is a step, not a slope.} At batch 32, one grammar costs
82\,\textmu{}s and two cost 388; but sixteen also cost 394. The 5$\times$ is paid
once, when the pool stops being homogeneous, and then does not grow with the
mixture. This is the property that matters for serving, where the mixture is
never one.

\textbf{We lose at two schemas and small batch.} 0.47$\times$ at batch 32 with
two schemas is the worst cell in the table, and it is the direct consequence of
the step above: we pay the full heterogeneity cost while XGrammar still has only
64 cheap host calls to make.

\textbf{The absolute cost at 8--16 schemas and batch 512 is bad.} 4\,ms is
33--36$\times$ what the same schemas cost run alone, against 0.82--1.04$\times$
for XGrammar. We still win the comparison 8.9--11.1$\times$ because their
absolute cost is 36--45\,ms, but the honest statement is that pool-wide ceilings
--- the replay window, the readings a group may have --- are sized by the
largest member and every sequence pays for it. This is the clearest target for
future work.

\subsection{Speculative decoding}
\label{sec:eval-spec}

A draft of $k$ tokens must be checked against the grammar and the accepted
prefix kept. This is the case where a host-side matcher stops being merely
costly. XGrammar offers \texttt{traverse\_draft\_tree}, which walks the whole
draft tree in one call and is close to flat in $k$; comparing against its
per-position path instead --- as an earlier version of this work did --- overstates
the gap by roughly 46$\times$ and we do not do so.

\begin{table}[t]
\caption{Speculative step, median \textmu{}s, against
\texttt{traverse\_draft\_tree}. Every position pays a fill and an advance on our
side, and the rollback is charged. Source: \texttt{probe\_speculative},
2026-07-31.}
\label{tab:spec}
\vskip 0.1in
\begin{center}
\begin{small}
\begin{tabular}{rrrrr}
\toprule
batch & $k$ & engrain & XGrammar & ratio \\
\midrule
\multirow{3}{*}{128}
 & 1 & 102 & 260 & 2.56$\times$ \\
 & 4 & 219 & 562 & 2.57$\times$ \\
 & 8 & 361 & 485 & 1.35$\times$ \\
\midrule
\multirow{3}{*}{512}
 & 1 & 278 & 1{,}072 & 3.85$\times$ \\
 & 4 & 474 & 1{,}999 & 4.22$\times$ \\
 & 8 & 730 & 1{,}980 & 2.71$\times$ \\
\bottomrule
\end{tabular}
\end{small}
\end{center}
\vskip -0.1in
\end{table}

We win at every $k$ measured (Table~\ref{tab:spec}), but the trend is against
us: our cost grows with $k$ (278/474/730 at batch 512) because each draft
position pays a fill and an advance, while theirs is roughly flat. Extrapolating,
the advantage closes near $k\approx16$. The structural point survives ---
our walk is one recorded graph replayed once, with rollback on device, and needs
no synchronization per position --- but the arithmetic advantage is finite and we
say so.

\subsection{The two strongest counter-arguments}
\label{sec:eval-counter}

Per-step numbers invite two objections, and both are stronger than any ratio we
can quote. We measured them rather than argued with them.

\paragraph{``XGrammar overlaps its fill with the forward pass, so residency buys
nothing.''} This is the objection that once forced this project to withdraw a
claim. We measure a real decode forward pass, each engine alone, and the two
together on one stream, and charge each engine only the time it \emph{adds} to
the step (Table~\ref{tab:overlap}). At batch 32 the objection holds: we add
44\,\textmu{}s and XGrammar adds 46, which is parity. It stops holding as the
batch grows, because what overlaps is bounded by the forward pass and
XGrammar's cost is not: at batch 512 it adds 398\,\textmu{}s to a
22\,ms step against our 90.

\begin{table}[t]
\caption{Cost \emph{added} to a real decode step when the constraint is
overlapped with the forward pass. The comparison the objection asks for.
Source: \texttt{engrain\_lab.rigor.overlap}.}
\label{tab:overlap}
\vskip 0.1in
\begin{center}
\begin{small}
\begin{tabular}{rrrrr}
\toprule
batch & forward & engrain adds & XGrammar adds & ratio \\
\midrule
32  & 6{,}559\,\textmu{}s & 44\,\textmu{}s & 46\,\textmu{}s  & 1.05$\times$ \\
128 & 9{,}394\,\textmu{}s & 60\,\textmu{}s & 112\,\textmu{}s & 1.87$\times$ \\
512 & 22{,}090\,\textmu{}s & 90\,\textmu{}s & 398\,\textmu{}s & 4.42$\times$ \\
\bottomrule
\end{tabular}
\end{small}
\end{center}
\vskip -0.1in
\end{table}

\paragraph{``The host is idle anyway.''} A host-side matcher is cheap when it
has a core to run on. Serving hosts do not offer that guarantee: the same CPUs
run detokenization, scheduling, request parsing and the API server. We sweep
the number of busy cores next to the fill (Table~\ref{tab:contention}). Our p50
and p99 are flat --- 27.8 to 27.9\,\textmu{}s and 51.0 to 56.7 --- because the
work is on the device. XGrammar's median degrades mildly (1.09$\times$) but its
\textbf{p99 goes from 449.6 to 3{,}495.3\,\textmu{}s, a 7.8$\times$ tail
blow-up}. Constrained decoding is a tail-latency business, and this is the
clearest measurement in the paper of what residency is actually for.

\begin{table}[t]
\caption{Mask fill at batch 128 while $n$ cores are kept busy beside it.
Source: \texttt{engrain\_lab.rigor.serving}.}
\label{tab:contention}
\vskip 0.1in
\begin{center}
\begin{small}
\begin{tabular}{rrrrr}
\toprule
busy & \multicolumn{2}{c}{engrain} & \multicolumn{2}{c}{XGrammar} \\
cores & p50 & p99 & p50 & p99 \\
\midrule
0  & 27.8 & 51.0 & 406.5 & 449.6 \\
8  & 27.5 & 66.5 & 433.2 & 464.2 \\
16 & 27.8 & 59.9 & 421.5 & 448.3 \\
24 & 27.9 & 56.7 & 444.7 & \textbf{3{,}495.3} \\
\bottomrule
\end{tabular}
\end{small}
\end{center}
\vskip -0.1in
\end{table}

\subsection{Ablation: what each mechanism is worth}
\label{sec:eval-ablation}

Table~\ref{tab:ablation} removes one mechanism at a time. Both knobs are set
\emph{before} the step is captured, which matters more than it sounds: our
first attempt set them afterwards and reported that the memo was worth
1.00$\times$ while hitting 90\% of steps --- the captured graph had the old
scalar arguments baked in, so nothing had actually been disabled. An ablation
of a graph-resident system has to be recorded, not configured.

\begin{table}[t]
\caption{Ablation on a three-schema batch with sequences at random document
positions. ``memo'' compares the full memo against a single-slot one; at batch
1 that is degenerate, since one slot always serves one sequence.
Source: \texttt{exp\_ablation}.}
\label{tab:ablation}
\vskip 0.1in
\begin{center}
\begin{small}
\begin{tabular}{rrrrr}
\toprule
batch & full & memo worth & capture worth & memo hit \\
\midrule
1   & 49.8\,\textmu{}s    & 1.00$\times$\,$^\dagger$ & 2.56$\times$ & 100.0\% \\
8   & 59.6\,\textmu{}s    & 6.15$\times$ & 2.34$\times$ & 100.0\% \\
32  & 408.6\,\textmu{}s   & 1.13$\times$ & 1.04$\times$ & 93.8\% \\
128 & 448.7\,\textmu{}s   & 1.16$\times$ & 1.04$\times$ & 90.6\% \\
512 & 3{,}980.4\,\textmu{}s & 1.00$\times$ & 1.00$\times$ & 86.1\% \\
\bottomrule
\end{tabular}
\end{small}
\end{center}
\vskip -0.1in
\end{table}

The two mechanisms pay in the same place and for the same reason: at small
batch the step is dispatch-bound and mostly redundant, so capture is worth
2.3--2.6$\times$ and the memo up to 6.2$\times$; at batch 512 both are
1.00$\times$ because the step is work-bound and the work is genuinely
different per sequence. We include the negative half of this result because a
mechanism that helps only at batch 8 should be described that way.

The third ablation is the fill split (\S\ref{sec:fill}), and it is the one that
justifies the paper's implementation lesson: on a balanced batch the single
fused kernel is fine, and on a skewed one it costs 9{,}883\,\textmu{}s against
691\,\textmu{}s. The regression is invisible to any benchmark that does not
seed sequences apart.



\begin{table}[t]
\caption{Constrained sampling at batch 512, measured in place (no restore of a
310\,MB logits tensor that a decode loop never pays). ``Shortlist'' emits
whichever of the allowed set and its complement is shorter, decided on device.
Source: session measurements, 2026-07-30.}
\label{tab:sampling}
\vskip 0.1in
\begin{center}
\begin{small}
\begin{tabular}{lrr}
\toprule
regime & vs.\ mask-then-sample & note \\
\midrule
all rows sparse           & 7.73$\times$ & upper bound \\
all rows dense            & 0.95$\times$ & compaction loses \\
distribution-weighted     & 1.79$\times$ & shortlist, per-step choice \\
ragged batch, 50\% dense  & 1.05$\times$ & \textbf{measured steady state} \\
ragged, sync-free scheme  & 0.86$\times$ & fixed shape, no host sync \\
\bottomrule
\end{tabular}
\end{small}
\end{center}
\vskip -0.1in
\end{table}

Table~\ref{tab:sampling} is the result we most want to be different. Fused
sampling is the cleanest consequence of residency --- constraining makes sampling
\emph{cheaper}, which is impossible if the constraint is a host product --- and on
a homogeneous sparse batch it delivers 7.7$\times$. But the measured steady
state of a continuously batched workload is about half dense rows
(\S\ref{sec:background}), and there the advantage is 1.05$\times$: parity.

We report the cause precisely because it is not in our engine. Two interfaces
prevent the ragged case from paying: XGrammar's
\texttt{apply\_token\_bitmask\_inplace} takes its row indices as a host
\texttt{List[int]}, and a sampler's output shape is its row count, so selecting a
subset of rows requires knowing that count on the host. What is needed is a
sampler that accepts a device-side row count. Until sampling APIs offer one, the
fused-sampling benefit is available in principle and 5\% in practice on a mixed
batch.

\subsection{Correctness, in two layers}
\label{sec:correctness}

Constrained decoding admits no approximation: one wrong bit either leaks an
illegal token or removes a legal one, so learned or lossy mask representations
are out of scope by construction. But ``correct'' names two different
properties, and conflating them is how this literature reports numbers that
cannot be acted on:

\begin{description}
\item[runtime exactness] the device mask equals what the compiled grammar
allows;
\item[lowering fidelity] the compiled grammar equals the schema.
\end{description}

The first is what residency depends on and we verify it exhaustively. The
second is where every grammar-based system loses ground, and it is where we
found --- and fixed --- a real defect in our own compiler.

\paragraph{Runtime exactness.} Every mask row the device produces is compared
bit-for-bit against an independent CPU reference matcher built from the same
tables. Four verifications run over the real corpus: \textbf{corpus} (619
rows over model-generated documents), \textbf{walk} (4{,}800 rows over random
walks, which reach states real documents do not), \textbf{conflicted} (687 rows
restricted to grammars with reduce/reduce conflicts, i.e.\ where forking is
live), and \textbf{mixed} (1{,}092 rows over batches of several grammars,
including a CUDA graph \emph{recorded on one assignment of grammars to
sequences and replayed on a different one}). All 7{,}198 rows match, with zero
failures. The graph-replay check is the one the residency claim rests on: a
recorded step survives the batch composition changing underneath it, which is
what continuous batching does every step. In addition, 99 unit tests run under
all three backends, and the \texttt{differential} backend runs the CUDA and
Triton implementations on every call and compares.

\paragraph{Lowering fidelity, and a metric that was measuring the wrong thing.}
The aggregate we had been reporting --- ``56.0\% of documents generated under
our mask validate, against XGrammar's 77.6\%'' --- turns out to conflate two
unrelated quantities, because a random walk that runs out of byte budget
produces truncated JSON and was counted as a failure. Decomposing it
(Table~\ref{tab:soundness}) separates \emph{completion} (did the walk finish a
document) from \emph{over-acceptance} (given a complete document, does the
schema accept it). Only the second is a property of the engine.

\begin{table}[t]
\caption{Random-walk generation under each engine's mask on the 194 schemas
both engines compile, same seed, same 400-byte budget. ``Before'' and ``after''
are the compiler defect described below. Source:
\texttt{exp\_soundness\_split}.}
\label{tab:soundness}
\vskip 0.1in
\begin{center}
\begin{small}
\begin{tabular}{lrrr}
\toprule
 & completion & valid $|$ complete & aggregate \\
\midrule
engrain, before & 52.5\% & 70.0\% & 36.7\% \\
engrain, after  & 50.5\% & \textbf{80.4\%} & 40.6\% \\
XGrammar        & 63.7\% & 94.6\% & 60.2\% \\
\bottomrule
\end{tabular}
\end{small}
\end{center}
\vskip -0.1in
\end{table}

\paragraph{Attributing the gap to a keyword found a bug.} Recording which JSON
Schema keyword rejected each over-accepted document showed the failure was not
diffuse: \textbf{219 of 243} over-acceptances were the single keyword
\texttt{required}, against 8 for XGrammar. That concentration is what made the
cause findable. It was not the precision level --- forcing the most faithful
lowering moved 181 of 194 schemas from \texttt{Shadowed} to \texttt{Exact},
cost 4$\times$ compile time and 2$\times$ tables, and left validity
\emph{unchanged} at 70.0\%, which ruled out the hypothesis we started with.

The cause was in the order-free object construction. To let properties arrive
in any order while still enforcing \texttt{required}, the compiler enumerates
subsets of the required set and carries the subset in the parser state, one
rule per subset. For the \emph{full} subset --- every required property seen ---
the rule that picks the next property has no alternatives when the object also
closes its property names (\texttt{additionalProperties: false}). The code
treated ``no property may follow'' as a construction failure and fell back to an
object that does not enforce \texttt{required} at all. The condition is not a
failure; it means the object must now close. Minimally: an object with one
declared property, that property required, and \texttt{additional\-Properties:
false} accepted \texttt{\{\}}. With two properties it did not, which is why no
unit test caught it.

Fixing it removes 78 of the 243 over-acceptances, eliminates the
\texttt{oneOf} failures entirely, and moves validity from 70.0\% to 80.4\%
(Table~\ref{tab:soundness}). We report the before and after because the
methodology --- attribute, minimize, fix, re-measure --- is the contribution
here, not the number.

\paragraph{What remains, and the honesty defect behind it.} 141
\texttt{required} failures remain, and they have a different cause: the subset
enumeration costs $2^{n}$ rules, so it is bounded at 4 required properties when
the property names are open and 6 when they are closed. Past that the object
widens. That is a defensible budget --- the alternative, fixing the property
order as XGrammar does, rejects valid permutations, which this system's
contract forbids --- but it was being applied \emph{silently}: 33 of the 34
corpus schemas whose \texttt{required} was dropped reported no approximation at
all, because the approximation list was derived from the precision level rather
than from what the lowering did. The engine now reports it in every case, which
takes the silent count to zero. We consider a constrained decoder that relaxes
a schema without saying so to be the more serious of the two defects, and we
suspect the distinction between coverage, exactness, and \emph{declared}
exactness is under-reported across this literature.

\subsection{The costs of residency}
\label{sec:costs}

Residency is not free, and the honest accounting is the following three
results. All are measured over 528 distinct JSONSchemaBench schemas.

\begin{table}[t]
\caption{What residency costs. Compile time and table memory are measured over
528 JSONSchemaBench schemas; XGrammar's figures are for the same schemas.
Source: \texttt{engrain\_lab.rigor.cost}, 2026-07-30.}
\label{tab:costs}
\vskip 0.1in
\begin{center}
\begin{small}
\begin{tabular}{lrrr}
\toprule
 & engrain & XGrammar & ratio \\
\midrule
compile p50      & 112\,ms   & 16\,ms  & 7.0$\times$ worse \\
compile p90      & 1.38\,s   & ---     & \\
compile p99      & 4.63\,s   & 0.6\,s  & 7.7$\times$ worse \\
compile max      & 12.5\,s   & ---     & \\
\midrule
tables p50       & 1.41\,MiB & 52\,KiB & 27$\times$ worse \\
tables p99       & 37.5\,MiB & ---     & 700$\times$ worse \\
tables max       & 97.6\,MiB & ---     & \\
\midrule
schemas compiled & 502/528 (95.1\%) & $\approx$98\% & $-$3\,pts \\
\bottomrule
\end{tabular}
\end{small}
\end{center}
\vskip -0.1in
\end{table}

\textbf{Compile time} (Table~\ref{tab:costs}) is 7$\times$ XGrammar's at the
median and worse in the tail. It is, however, \emph{linear} in the number of
lexer states (slope 1.05, $r=0.94$): there is no algorithmic blow-up, only the
inherent $O(|V| \times \text{lexer states})$ cost of computing what every token
does from every lexer state. In a serving context, schemas arrive with requests,
so this is a stated failure and not a footnote.

\textbf{Memory.} 1.41\,MiB at the median against a 52\,KiB compiler cache. We
looked for structural sharing and did not find much: whole-array sharing across
schemas gives 1.0$\times$ (nothing is byte-equal), group sets are already
interned, and interning reading lists saves 12\%. The one lever with the right
shape is building a lexer state's reading tables \emph{on demand}, the first
time a parse reaches that state, since a real document visits only 1--4\% of the
states its grammar can reach. That would address compile time and memory with
one mechanism, and it is the main piece of future work.

\textbf{Coverage.} 502 of 528 schemas compile (95.1\%). The 26 refusals are
typed: 12 the front end cannot lower, 10 exceed the production budget, 4 exceed
the lexer budget. Separately --- and this is a distinction the constrained
decoding literature does not usually draw --- 371 of the 502 compiled grammars
(73.9\%) accept a \emph{relaxation} of their schema rather than exactly the
schema. The relaxations come from the lowering, not the runtime: \texttt{oneOf}
becomes a union, and a union is ``at least one'' where \texttt{oneOf} is
``exactly one''. The runtime only ever narrows. We report bit-exactness against
\emph{our compiled grammar}, which is a weaker and more precise claim than
bit-exactness against the schema, and we believe every system in this space owes
the same distinction.

\section{Limitations}
\label{sec:limitations}

We state these as results rather than as caveats, because the design's central
claim is structural and survives them, while a reader deciding whether to deploy
needs them.

\begin{enumerate}
\item \textbf{Batch 1 is a loss} (0.24--0.90$\times$). Our floor is set by
sweeping live configurations, not by the batch.
\item \textbf{Compile is 7$\times$ slower} and the p99 is seconds. Lazy lexer
residency is the plausible fix and is unimplemented.
\item \textbf{Tables are 27--700$\times$ larger} than a host matcher's cache.
\item \textbf{Our lowering over-accepts 19.6\% of complete documents}, against
XGrammar's 5.4\%, and 141 of the 153 remaining cases are one keyword,
\texttt{required}, dropped where the subset enumeration exceeds its budget.
\item \textbf{Fused sampling is 1.05$\times$ on a ragged batch}, not the
7.7$\times$ a homogeneous one suggests, and the blocker is in surrounding APIs.
\item \textbf{Heterogeneous pools of 8--16 schemas cost us 33--36$\times$}
what those schemas cost alone.
\item \textbf{The memo and graph capture pay only below batch 32}
(Table~\ref{tab:ablation}); at batch 512 both are 1.00$\times$.
\item \textbf{One baseline, one GPU.} We compare against XGrammar; llguidance
is installed and not yet in the comparison. Every launch shape, memo size and
chunk count was fitted on an A100.
\item \textbf{No end-to-end throughput claim.} Grammar work is a few percent of
a decode step, and a vLLM A/B at batch 256 ranged from 4{,}523 to 13{,}325
tokens per second across runs --- too noisy to attribute a difference. We report
the per-step result, which is the defensible one.
\end{enumerate}

\section{Related Work}
\label{sec:related}

\paragraph{Host-side constrained decoding.} Outlines~\cite{willard2023efficient}
indexes a regular-language DFA against the vocabulary;
SynCode~\cite{ugare2025syncode} and grammar-aligned
decoding~\cite{park2024grammaraligned} extend to context-free grammars;
XGrammar~\cite{dong2024xgrammar} makes the pushdown automaton practical with
adaptive token masks and is the deployed baseline;
llguidance and Domino~\cite{beurerkellner2024domino} attack the same overhead
with better host algorithms and speculation. All compute the mask on the CPU,
which is the property this paper is about.

\paragraph{GPU-side work.} Pre3~\cite{chen2025pre3} and PSC~\cite{li2026psc}
precompute deterministic pushdown structure to make the host step cheaper or to
classify parser stacks; Gram2Token~\cite{hua2026gram2token} moves the
token-to-terminal bridge toward the device. These reduce the cost of the host
step or move part of it; to our knowledge none keeps the parser stack itself on
the device across steps, which is what makes graph capture and sampler fusion
possible.

\paragraph{Automata and tokenizers.} The tokenizer--grammar impedance mismatch
is treated by~\cite{koo2024automata} and~\cite{park2025flexible};
CRANE~\cite{banerjee2025crane} studies the expressiveness cost of constraining.
Our token-group quotient is a device-oriented instance of the same idea.

\paragraph{Serving.} CUDA graph capture of the decode loop is standard in
vLLM and SGLang~\cite{zheng2024sglang}; our contribution is to make the grammar
a participant in it.

\section{Conclusion}

Constrained decoding has been treated as a host-side algorithms problem, and by
that measure it is nearly solved. We argued that the interesting question is
where the parser's state lives, because a decode step that a serving engine
records as a CUDA graph can only contain device work. Engrain keeps an LALR(1)
parser --- lexer state, stack, and conflict forks --- resident on the GPU, which
makes the constraint capturable, fusable with the sampler, and advanceable
through a speculative draft without a host round trip. It costs 103--327\,\textmu{}s
per full step at batch 512 against 2{,}041--6{,}811 for the deployed baseline,
1.35--4.22$\times$ less under speculation, and is bit-exact over 7{,}198 verified
rows. It also costs 7$\times$ more to compile and 27$\times$ more memory, loses
below batch 16, and delivers only parity on fused sampling once the batch is
ragged. We believe the structural claim is worth those prices, and we have tried
to report them precisely enough that a reader can disagree.

\bibliography{references}
\bibliographystyle{mlsys2025}

\end{document}
